% Practical C++ reference and guided exercises. The separate problem and
% pattern indexes were removed; the redesigned contents page is the book's
% single navigation system.

\clearpage
\phantomsection
\section*{C++ Reference and Guided Practice}
\addcontentsline{toc}{section}{C++ Reference and Guided Practice}
\markboth{C++ Reference and Guided Practice}{Reference}

The solutions use the signatures expected by an online judge. This reference
collects the node definitions that are normally supplied by the judge and
clarifies the difference between borrowed links and owned objects. It also
provides a small set of exercises for turning the completed lessons into
active interview practice.

\begin{conceptbox}{Implementation validation}
All 75 complete solutions are checked independently with a C++
compiler. The audit includes judge-provided tree, list, and graph types, so a
solution cannot pass merely because another chapter supplied a missing type or
header. Representative examples and boundary cases are also exercised for
each algorithm family.
\end{conceptbox}

\subsection{Judge-provided pointer structures}

\begin{visualbox}{Values and links in the three common node types}
\centering
\begin{tikzpicture}[font=\footnotesize,>={Stealth[length=2.2mm,width=1.55mm]}]
  \node[dsnode,active,minimum width=24mm] (tree) at (0,1.5)
    {\texttt{TreeNode}\\value};
  \node[dsnode,minimum width=19mm] (tl) at (-2.0,0) {left child};
  \node[dsnode,minimum width=19mm] (tr) at (2.0,0) {right child};
  \draw[flowarrow] (tree.south west) -- (tl.north);
  \draw[flowarrow] (tree.south east) -- (tr.north);

  \node[dsnode,active,minimum width=24mm] (list) at (5.2,1.5)
    {\texttt{ListNode}\\value};
  \node[dsnode,minimum width=19mm] (next) at (5.2,0) {next node};
  \draw[flowarrow] (list.south) -- (next.north);

  \node[dsnode,active,minimum width=24mm] (graph) at (2.6,-1.8)
    {graph node\\value};
  \node[dsnode,minimum width=22mm] (neighbors) at (2.6,-3.2)
    {neighbor pointers};
  \draw[flowarrow] (graph.south) -- (neighbors.north);
\end{tikzpicture}

\smallskip
The pointer fields describe structure. In judge code, the judge normally owns
the input nodes and the solution temporarily follows or rearranges those links.
\end{visualbox}

\subsection{Common judge definitions}

\begin{lstlisting}[style=compactcode,float=htbp,caption={Typical linked-list and binary-tree definitions}]
using namespace std;

struct ListNode {
    int val;
    ListNode* next;
    ListNode(int x) : val(x), next(nullptr) {}
};

struct TreeNode {
    int val;
    TreeNode* left;
    TreeNode* right;
    TreeNode(int x) : val(x), left(nullptr), right(nullptr) {}
};
\end{lstlisting}

\begin{lstlisting}[style=compactcode,float=htbp,caption={Typical graph-node definition}]
#include <utility>
#include <vector>

using namespace std;

class Node {
public:
    int val;
    vector<Node*> neighbors;

    Node() : val(0) {}
    Node(int value) : val(value) {}
    Node(int value, vector<Node*> links)
        : val(value), neighbors(move(links)) {}
};
\end{lstlisting}

\begin{conceptbox}{Raw pointers and ownership}
The examples use raw pointers because the judge owns the supplied structures
and requires standard method signatures. Production C++ should make ownership
explicit with RAII containers and smart pointers where one object genuinely
owns another. A map used while cloning a graph stores lookup pointers; it does
not become a second owner of the original or cloned nodes.
\end{conceptbox}

\subsection{A reliable code-review pass}

Before submitting a solution, read it once without thinking about the example.
Check the following properties directly in the code.

\begin{revisionbox}{Five checks before submission}
\begin{enumerate}
  \item Every input permitted by the rules reaches a valid base case or
  loop termination condition.
  \item Array and string indexes are checked before access, especially after a
  pointer moves or a window shrinks.
  \item Values needed by an update are copied before the original variables are
  overwritten.
  \item Recursive and graph solutions mark, unmark, or memoize state at the
  exact point required by their key rule.
  \item The stated complexity counts sorting, heap operations, recursion depth,
  and output storage consistently.
\end{enumerate}
\end{revisionbox}

\subsection{Guided practice with hints}

For each exercise, first draw the changing state, then write the key rule and
only afterward modify the implementation.

\begin{revisionbox}{Practice prompts}
\small
\renewcommand{\arraystretch}{1.08}
\begin{tabularx}{\linewidth}{@{}>{\raggedright\arraybackslash\bfseries}p{2.45cm}
  >{\raggedright\arraybackslash}X
  >{\raggedright\arraybackslash}X@{}}
\toprule
Area & Exercise & Hint after drawing \\
\midrule
Arrays and greedy & Return the buy and sell days, not only the profit. & Save
the day when the running minimum changes and snapshot both days when the best
profit improves. \\
Two pointers & Generalize 3Sum to an arbitrary target. & Compare the total to
the requested target; the safe pointer directions do not change. \\
Trie recursion & Count every word matching a wildcard pattern. & Sum successful
branches instead of returning after the first match. \\
Tree queues & Produce zigzag level order. & Keep the BFS queue unchanged and
alter only the destination index inside each row. \\
Graphs & Return one valid course order. & Record vertices when they leave a
zero-indegree queue; a short result still detects a cycle. \\
Dynamic programming & Reconstruct a minimum coin selection. & Store the final
coin responsible for each improving transition. \\
Heaps & Support deletion from a median stream. & Combine the heaps with delayed
deletion counts in hash maps. \\
Strings & Decode payloads containing digits and \texttt{\#}. & The stored
length, not a delimiter search, controls how many payload characters to read. \\
\bottomrule
\end{tabularx}
\end{revisionbox}